% ===================== ZUSAMMENFASSUNG (ANCHO COMPLETO) =====================

\begin{tabular}{p{\linewidth}}
  \footnotesize
  \justifying
    \cellcolor{cvbackgroundgray}
    \noindent
Masterstudent der Computational Science and Engineering an der Universität Ulm, Machine Learning Engineer und Maschinenbauingenieur mit Erfahrung an der Schnittstelle von Big Data, Künstlicher Intelligenz und technischen Entwicklungssystemen. Schwerpunkt meiner Arbeit ist die datengetriebene Analyse industrieller Signale sowie die Entwicklung robuster KI-Modelle für Signalvorhersage, Anomalieerkennung und Entscheidungsunterstützung. Dabei arbeite ich mit Python, SQL, Azure ML, Databricks und skalierbaren ETL-Workflows zur Verarbeitung hochfrequenter Zeitreihendaten aus industriellen Umgebungen. Ich kombiniere fundierte Kenntnisse in Data Engineering, Machine Learning und mechanischen Systemen, um strukturierte, belastbare und praxisnahe KI-Lösungen für moderne Engineering- und Entwicklungsprozesse umzusetzen.
\end{tabular}



%
% hendsoldt dataenauswertung
% \noindent Machine-Learning-Ingenieur mit Erfahrung in der datenbasierten Analyse technischer Systeme und industrieller Prozesse. Fundierte Kenntnisse in Python-gestützter Datenauswertung, Datenaufbereitung und der Integration heterogener Datenquellen.

% Erfahrung in der Entwicklung strukturierter Analyse-Workflows von der Datenextraktion über die Analyse bis zur Visualisierung von Ergebnissen in Dashboards und Reports. Besonders interessiert an Aufgaben an der Schnittstelle zwischen technischer Datenauswertung, Fehleranalyse und praxisnahen Softwarelösungen für ingenieurtechnische Anwendungen.


% Data Engineer / ML mit Data-Engineering-Fokus

% Machine Learning Engineer mit starkem Data-Engineering-Schwerpunkt, erfahren in der Konzeption skalierbarer Datenarchitekturen und der Automatisierung hochfrequenter industrieller Zeitreihen-Pipelines für KI-gestützte Predictive-Maintenance-Systeme.
% Kompetent im Aufbau robuster ETL-Workflows, der Integration von Shopfloor-IoT-Datenströmen in Cloud-basierte Analytics-Plattformen (Azure ML, Snowflake, Databricks, OPC UA, MQTT) sowie in der zuverlässigen End-to-End-Implementierung von ML-Lösungen.
% Nachweisbarer wirtschaftlicher Mehrwert durch messbare Verbesserungen in Anlagenverfügbarkeit, Produktionsleistung und operativer Effizienz.
% Maschinenbau-Hintergrund zur Unterstützung physikbasierter Systemauslegung.

% Machine Learning

% Machine Learning Engineer und Maschinenbauingenieur, spezialisiert auf KI-gestützte Predictive Maintenance und physik-informierte Zustandsüberwachung industrieller Systeme.
% Erfahrung in Entwicklung, Validierung und Deployment von ML/DL-Modellen zur Anomalieerkennung und RUL-Prognose, unterstützt durch skalierbare Data-Engineering-Pipelines für hochfrequente Zeitreihendaten.
% Fundierte Integration von IoT-Datenströmen zwischen Shopfloor- und Cloud-Analytics-Ebenen (Azure ML, Snowflake, Databricks, OPC UA, MQTT).
% Nachweislicher Geschäftsnutzen durch messbare Verbesserungen in Anlagenverfügbarkeit und Produktionsperformance.

% Research-Oriented Profil

% Maschinenbauingenieur und Machine-Learning-Spezialist mit Fokus auf physik-informierter Modellierung und KI-basierter Zustandsüberwachung industrieller cyber-physischer Systeme.
% Ausgeprägte Fähigkeit zur Interpretation von Sensordaten, Analyse physikalischer Degradationsmechanismen sowie Definition physikalisch konsistenter Anomaliekriterien vor der Implementierung skalierbarer ML-Lösungen.
% Erfahrung in der Modellierung hochfrequenter industrieller Zeitreihen, Automatisierung von Datenpipelines und End-to-End-Systemdeployment.
% Nachweisbare industrielle Wirkung bei gleichzeitiger Verfolgung mathematisch fundierter Ansätze in Predictive Maintenance und strukturierter Systemmodellierung.

% Meine jüngste Arbeit legt zunehmend einen Schwerpunkt auf strukturierte und topologie-bewusste Systemrepräsentationen, die physikalische Beziehungen und Randbedingungen explizit kodieren. Dies spiegelt eine wachsende forschungsorientierte Ausrichtung hin zu interpretierbaren und physikalisch fundierten KI-Modellen für komplexe ingenieurtechnische Architekturen wider.